\documentclass[12pt, a4paper]{article}
\usepackage[utf8]{inputenc}
\usepackage[spanish]{babel}
\usepackage{amsmath, amssymb}
\usepackage{geometry}
\usepackage{enumitem}
\geometry{top=2.5cm, bottom=2.5cm, left=2.5cm, right=2.5cm}

\title{Resolución de Ejercicio: Técnicas de Conteo}
\author{}
\date{}

\begin{document}

\maketitle

\section*{Enunciado}
19. Se quiere elegir un comité formado por tres miembros, presidente, tesorero y secretario. Para seleccionarlo disponemos de cuatro candidatos: Arturo, Basilio, Carlos y David. ¿Cuántos comités diferentes se pueden elegir entre los cuatro candidatos? Ejemplo: que Arturo sea presidente, Carlos sea tesorero y David sea secretario.

\vspace{0.5cm}
\hrule
\vspace{0.5cm}

\section*{Solución Paso a Paso}

\subsection*{1. Análisis e identificación del modelo}
Para resolver el problema, evaluamos las características de la selección a realizar mediante tres preguntas clave:
\begin{itemize}[label=\checkmark]
    \item \textbf{¿Importa el orden?} Sí. Los roles dentro del comité son distintos (presidente, tesorero y secretario). No es lo mismo que Arturo sea presidente y Carlos tesorero, a que Carlos sea presidente y Arturo tesorero[cite: 1].
    \item \textbf{¿Intervienen todos los elementos?} No. Disponemos de $m=4$ candidatos totales, pero solo seleccionaremos un subconjunto de $n=3$ de ellos para formar el comité[cite: 1].
    \item \textbf{¿Se pueden repetir los elementos?} No. Una misma persona no puede ocupar más de un cargo al mismo tiempo en el comité (muestreo sin reposición)[cite: 1].
\end{itemize}

Dado que importa el orden, no se seleccionan todos los elementos y no hay repetición, el modelo combinatorio adecuado es el de \textbf{Variaciones sin repetición}[cite: 1].

\subsection*{2. Planteo matemático y cálculo}
La fórmula estadística para calcular los subconjuntos ordenados de $n$ elementos que pueden seleccionarse de un conjunto de $m$ elementos es[cite: 1]:
\[ V_{(m,n)} = \frac{m!}{(m-n)!} \]

Definimos nuestra población total $m = 4$ (candidatos) y el tamaño del subconjunto $n = 3$ (cargos), y desarrollamos:
\begin{align*}
V_{(4,3)} &= \frac{4!}{(4-3)!} \\
V_{(4,3)} &= \frac{4!}{1!} \\
V_{(4,3)} &= \frac{4 \times 3 \times 2 \times 1}{1} \\
V_{(4,3)} &= 24
\end{align*}

Alternativamente, aplicando la regla de la multiplicación: hay 4 opciones para el cargo de presidente, 3 opciones restantes para el de tesorero, y 2 opciones finales para el de secretario ($4 \times 3 \times 2 = 24$)[cite: 1].

\subsection*{3. Conclusión}
Existen exactamente 24 formas diferentes de asignar los tres cargos jerárquicos entre los cuatro candidatos disponibles.

\vspace{0.5cm}

\textbf{Respuesta:} Se pueden elegir \textbf{24 comités} diferentes.

\end{document}
